\section{Charges for which the potential is not monotone, and questions}\label{sec:remarks}

The phase potential is monotone for the charges of
Theorem~\ref{thm:monotone}. This section gives charges for which it is not,
and states the questions left open.

\subsection{Reversed and ordered charges}
\begin{proposition}\label{prop:reversed}
If all $A_j=0$, then for all charges in the upper half-plane
\[
 \cD=c_0\bigl[\,n_0(\alpha_0-\alpha_1)+\pi\tr e^{-L}\,\bigr].
\]
In particular, for reversed charges one has $\cD<0$ at every such state with
$\pi\tr e^{-L}>n_0(\alpha_1-\alpha_0)$; such states exist if and only if
$\pi r>n_0(\alpha_1-\alpha_0)$.
\end{proposition}

\begin{proof}
The matrices $\Omega_V=z_0e^{-B}$, $-\Omega_U=z_0e^{-L}$, $\Omega_W=z_1\id$
are normal with phases $\alpha_0$, $\alpha_0-\pi$, $\alpha_1$. Hence
$\cD_V=|z_0|(n_0+\tr e^{-L})(\alpha_0-\varphi)\sin(\alpha_0-\varphi)$,
$\cD_U=|z_0|\tr e^{-L}(\pi-\alpha_0+\varphi)\sin(\alpha_0-\varphi)$ and
$\cD_W=n_1c_1(\varphi-\alpha_1)=n_0c_0(\varphi-\alpha_1)$. For reversed
charges $c_0<0$, and $\tr e^{-L}$ takes all values in $(0,r)$.
\end{proof}

For example, for $z_0=1+2\ii$, $z_1=-3+2\ii$, $n_0=n_1=r=1$, $D=(0,\frac12)$
and standard metrics, $\cD=-1.78927\ldots$

For ordered charges the potential can also fail to be monotone when
$r\geq2$.

\begin{proposition}\label{prop:orderedfailure}
Let $z_0=10^{-6}(-1+\ii)$ and $z_1=1+\frac\ii{500}$, so that
$\alpha_0=\frac{3\pi}4>\alpha_1$, and let $n_0=n_1=1$, $r=2$, $k=2$. With
standard metrics on $V=\CC^3$, $U=\CC^2$, $W=\CC$, let
\[
 D=\begin{pmatrix}0&\sqrt{\ell_1}&0\\0&0&\sqrt{\ell_2}\end{pmatrix},\quad
 A_1=(\tfrac15,0,0),\quad A_2=(0,a_1,a_2),\qquad\ell_1=\tfrac1{20},\ \ell_2=8,
\]
\[
 a_1=\sqrt{\epsilon/\ell_1},\quad a_2=\tfrac7{10}\sqrt{\epsilon\,\ell_2},\quad
 \epsilon=\tfrac3{10}\Bigl(\frac{1-e^{-\ell_1}}{\ell_1^2}+\frac{49}{100}(1-e^{-\ell_2})\Bigr)^{-1}.
\]
The cohomology representation is the nonzero map $\frac15:\CC\to\CC$, which
is stable, and
\[
 \cD=-1.6654354\ldots\times10^{-6}<0,\qquad
 (\cD_V,\cD_U,\cD_W)\approx(6.72,\,-9.09,\,0.70)\times10^{-6}.
\]
\end{proposition}

\begin{proof}
Since $C=\diag(\frac1{25},\,(a_1,a_2)^\dagger(a_1,a_2))$ has no block between
$\ker D$ and its complement, all blocks of $\Omega$ are of size $1$ or
$2$, and the polar decomposition of a $2\times2$ matrix is explicit, as in
the proof of Proposition~\ref{prop:mass}. The value of $\cD$ was enclosed
with interval arithmetic at $300$ bits; the enclosure has radius below
$10^{-30}$.\footnote{The program
\texttt{prop\_10\_1\_10\_2\_dissipation.py} and its output accompany this
paper as supplementary files.} At the same state and for the symmetric charges the same
computation gives $\cD=3.79\ldots>0$, in agreement with
Theorem~\ref{thm:monotone}.
\end{proof}

The mechanism is the following. As $z_0\to0$ with $z_1=e^{\ii\epsilon'}$ and
no arrows on the cohomology, $\Omega_V$ and $\Omega_U$ tend to the Hermitian
matrices $K_{F^\dagger F}(C_T)$ and $J$, whose phases have the eigenvalues
$0,\pi$ and $0,-\pi$ respectively, and
$\cD/\sigma\to\epsilon'\,\tr Y/(1-\tr Y)+\pi\bigl(\tr K_{F^\dagger F}(C_T)_--\tr J_-\bigr)$
up to terms of order $\epsilon'^2$ ($n_1=1$, $\tr Y<1$). Configurations with
$\tr J_->\tr K_{F^\dagger F}(C_T)_-$ exist, and then $\cD<0$ for small
$\epsilon'$. By Remark~\ref{rem:sharp} the block of degree one alone is
negative on a much larger set of charges; in numerical minimizations of
$\cD$ over the maps $D,A_j$, which are not used in any proof, the blocks of
degree zero compensated it except for charges very close to the real axis. By
Theorems~\ref{thm:cancellation} and~\ref{thm:convergence} the failure of
monotonicity of the potential does not affect the long-time behaviour.

\subsection{Questions}
We collect the questions left open.
\begin{enumerate}
\item Determine the best constant in Theorem~\ref{thm:kernel}. The value
 $(\sqrt2-1)/2$ would enlarge the admissible range to
 $\varphi\geq\frac\pi2(\sqrt2-1)$ and improve the constant $\rho$ of
 Theorem~\ref{thm:energy}.
\item Determine the set of charges for which $\cD\geq0$ at every state. It
 contains the admissible charges, and by
 Proposition~\ref{prop:orderedfailure} it does not contain all ordered
 charges when $r\geq2$.
\item For semistable cohomology that is not polystable, does the bounded
 factor between the induced metrics and the asymptotic solutions of
 \cite{hkkp23} converge, as in Proposition~\ref{prop:semistable}
 (Remark~\ref{rem:semistable}; \cite[Corollary~6.5]{wue26b})?
\item Is the mass nonincreasing along the flow for the symmetric charges and
 few arrows?
\item Do the parameter derivatives of $\cK_t$ converge pointwise on
 $\mathcal P^{\mathrm s}$ (Theorem~\ref{thm:families})?
\item Extend Theorems~\ref{thm:cancellation} and~\ref{thm:convergence} to
 complexes with acyclic parts at both vertices or in several degrees. The
 block of degree one then involves a Schur multiplier given by divided
 differences of the exponential function with one more interpolation point,
 to which Theorem~\ref{thm:kernel} does not apply.
\end{enumerate}
